\chapter{Resultados}
\label{cap:resultados}

Apresente os dados obtidos de forma organizada, sem interpreta-los ainda.
Tabelas e graficos carregam o peso deste capitulo.

\begin{table}[htb]
  \centering
  \caption{Sintese dos resultados obtidos}
  \label{tab:resultados}
  \begin{tabular}{lrrr}
    \toprule
    \textbf{Indicador} & \textbf{Previsto} & \textbf{Realizado} & \textbf{Desvio} \\
    \midrule
    Primeiro indicador & 100 & 112 & +12\% \\
    Segundo indicador  & 250 & 238 & $-$4,8\% \\
    Terceiro indicador &  80 &  80 & 0\% \\
    \bottomrule
  \end{tabular}
  \fonte{Elaborado pelo autor (2026).}
\end{table}

\section{Analise}
\label{sec:analise}

Agora sim, o que os numeros da Tabela~\ref{tab:resultados} significam: o que
saiu como esperado, o que destoou e por que.
